\section{Omitted Proofs}
\label{app:omitted_proofs}

This appendix collects the proofs omitted from the main body of the paper.

\subsection{Proof of Proposition~\ref{prop:mediumrun_dp}}

\begin{proof}
In the medium run a single worker performs the entire workflow, so the firm chooses only the AI strategy and pays the wage bill $\bigl(\sum_{T_b \in \T} \skillcost{b}\bigr)\bigl(\sum_{T_b \in \T} \timecost{b}\bigr)$, the product of the worker's total skill and total time.
Because the wage bill multiplies these two totals, the running value of one alone never reveals how an added task changes it, so the dynamic program carries both and optimizes the worker's job one task at a time.
For the proof we first show that the firm's medium-run optimization problem can be stated as a recursive problem.
Then proceed to use dynamic programming to approximate the optimal solution.

\emph{Recursive formulation.}
For $0 \le i \le m$ and skill and time levels $\skillcost{}$ and $\timecost{}$, let $R(i,\, \skillcost{},\, \timecost{})$ denote the minimum wage bill of a strategy for the whole workflow in which the worker has already been assigned a set of tasks, over steps $i+1, \dotsc, m$, with total skill $\skillcost{}$ and total time $\timecost{}$.
Here step $i$ is the highest-indexed step not yet assigned, so $R(i,\, \skillcost{},\, \timecost{})$ optimizes over all ways of covering steps $1, \dotsc, i$ and optimizing them within the same job.

In the base case $i = 0$ no steps remain to assign, the worker's totals are exactly $(\skillcost{}, \timecost{})$, and the wage bill is their product,
\[
R(0,\, \skillcost{},\, \timecost{}) = \skillcost{}\,\timecost{}, \qquad \text{for all } \skillcost{}, \timecost{}.
\]
The medium-run optimum is $R(m,\, 0,\, 0)$, the cheapest way to execute all $m$ steps within a job that starts with no tasks.

For $i \ge 1$ the firm chooses how step $i$ is assigned optimally, with $R(i,\, \skillcost{},\, \timecost{})$ taking the form
\begin{align*}
R(i,\, \skillcost{},\, \timecost{}) = \min\Biggl\{
& \ \underbrace{R\bigl(i-1,\; \skillcost{} + \manualSkill{i},\; \timecost{} + \manualTime{i}\bigr)}_{\text{Step $i$ is manual}}, \ \underbrace{\min_{0 \le r < i}\ R\!\left(r,\; \skillcost{} + \AIskill{i},\; \timecost{} + \frac{\AItime{i}}{\prod_{s=r+1}^{i} q_s}\right)}_{\text{Step $i$ optimally chained}}
\Biggr\}.
\end{align*}
In words, when assigning step $i$, the two options the firm considers are:
\begin{itemize}
\item \textbf{Manual step.}
Step $i$ is performed manually as its own task, contributing skill $\manualSkill{i}$ and time $\manualTime{i}$ to the cost and leaving steps $1, \dotsc, i-1$ to be assigned optimally:
\[
R\bigl(i-1,\; \skillcost{} + \manualSkill{i},\; \timecost{} + \manualTime{i}\bigr).
\]
\item \textbf{AI chain.}
Step $i$ is the augmented end of an AI chain over steps $r+1, \dotsc, i$ for some $r < i$, which automates steps $r+1, \dotsc, i-1$ and adds skill $\AIskill{i}$ and time $\AItime{i}/\prod_{s=r+1}^{i} q_s$ to the job, leaving steps $1, \dotsc, r$ to assign:
\[
\min_{0 \le r < i}\ R\!\left(r,\; \skillcost{} + \AIskill{i},\; \timecost{} + \frac{\AItime{i}}{\prod_{s=r+1}^{i} q_s}\right).
\]
The cut $r$ is the highest-indexed step left out of the chain.
The special case where $r = i-1$ gives a length-one chain that merely augments step $i$.
\end{itemize}


\emph{Discretization and Computation Time.}
Given our recursive formulation, we wish to use dynamic programming to solve for the optimal AI strategy by filling in the values of $R(i,\; \skillcost{},\; \timecost{})$ for all $i$, $\skillcost{}$, and $\timecost{}$.
Since $\skillcost{}$ and $\timecost{}$ take on continuous values, we will discretize all skill costs $\skillcost{}$ and time costs $\timecost{}$ to appropriate powers of $(1+\epsilon)$, where $\epsilon > 0$ is an arbitrarily small error term.

Suppose that all manual and augmented skill and time costs lie in $[1/B, B]$ for some $B > 0$.
Performing every step manually is feasible at a wage bill of at most $(mB)(mB) = m^2 B^2$, so no strategy with a larger wage bill need be considered.
A worker's total skill combines additively and lies in $[1/B,\, mB]$.
Since this skill is at least $1/B$, and since no strategy with wage bill above $m^2 B^2$ need be considered, the total time of any strategy worth considering is at most $m^2 B^2/(1/B) = m^2 B^3$, so time cost of optimal strategy lies in $[1/B,\, m^2 B^3]$.


We conclude that when filling table $R(i, \skillcost{}, \timecost{})$, it suffices to consider skill costs $\skillcost{}$ lying in the range $[1/B, mB]$ and time costs $\timecost{}$ lying in the range $[1/B, m^2 B^3]$.
Each of these two ranges contains $O\bigl(\epsilon^{-1}\log(mB)\bigr)$ powers of $(1+\epsilon)$.
The table $R(i,\, \skillcost{},\, \timecost{})$ is indexed by the step $i$, which takes $m+1$ values, together with the discretized skill and time levels $\skillcost{}$ and $\timecost{}$, each taking $O\bigl(\epsilon^{-1}\log(mB)\bigr)$ values, so it holds
\[
(m+1)\,\cdot\,O\bigl(\epsilon^{-1}\log(mB)\bigr)\,\cdot\,O\bigl(\epsilon^{-1}\log(mB)\bigr) \;=\; O\bigl(m\,\epsilon^{-2}\log^2(mB)\bigr)
\]
entries in total.
Each entry is filled in $O(m)$ time by taking the minimum over the manual option and the at most $m$ chain cuts $r$ in the recursion above, so filling the whole table takes
\[
O\bigl(m\,\epsilon^{-2}\log^2(mB)\bigr)\,\cdot\,O(m) \;=\; O\bigl(m^2\,\epsilon^{-2}\log^2(mB)\bigr)
\]
time.
Recording which option attains each minimum lets the firm read the optimal AI strategy off $R(m,\, 0,\, 0)$.

It remains to bound the discretization error.
The wage bill is evaluated only at the base case, where it multiplies the worker's total skill and total time; rounding each of these down to a power of $(1+\epsilon)$ moves it by a factor of at most $(1+\epsilon)$, so the product is computed to within $(1+\epsilon)^2$.
Along the recursion the manual and chain options accumulate skill and time additively, and rounding the running totals to powers of $(1+\epsilon)$ keeps a multiplicative error of at most $(1+\epsilon)$ in each.
Rescaling $\epsilon$ by a constant therefore yields the stated $(1+\epsilon)$ approximation.
\end{proof}

\subsection{Proof of Lemma~\ref{lem:ca_predecessor}}

\begin{proof}
The proof is fairly mechanical.
The two strategies agree on every step outside the block $k-1, k, k+1$, and both reach the endpoint $k+1$ as an augmented step, paying $\AIskill{k+1}\AItime{k+1}/q_{k+1}$ for its management.
Within the block, the described strategy costs
\[
\underbrace{\min\!\left\{\manualSkill{k-1}\manualTime{k-1},\ \frac{\AIskill{k-1}\AItime{k-1}}{q_{k-1}}\right\} + \manualSkill{k}\manualTime{k} + \frac{\AIskill{k+1}\AItime{k+1}}{q_{k+1}}}_{k-1\text{ on its own},\ k\text{ manual},\ k+1\text{ augmented}},
\]
while running the three steps as one AI chain augmented at $k+1$ costs $\AIskill{k+1}\AItime{k+1}/(q_{k-1}q_k q_{k+1})$, paying nothing to manage the predecessor or step $k$.
The chain is cheaper exactly when the \emph{marginal} cost of folding the pair $\{k-1, k\}$ in, the difference of the two chain terms, falls below the cost of doing that pair the old way.
Crediting back the common endpoint term $\AIskill{k+1}\AItime{k+1}/q_{k+1}$ from both sides,
\[
\frac{\AIskill{k+1}\AItime{k+1}}{q_{k-1}q_k q_{k+1}} - \frac{\AIskill{k+1}\AItime{k+1}}{q_{k+1}}
\;=\;
\frac{\AIskill{k+1}\AItime{k+1}\,(1-q_{k-1}q_k)}{q_{k-1}q_k q_{k+1}}
\;<\;
\min\!\left\{\manualSkill{k-1}\manualTime{k-1},\ \frac{\AIskill{k-1}\AItime{k-1}}{q_{k-1}}\right\} + \manualSkill{k}\manualTime{k},
\]
which is exactly condition~\eqref{eq:pred_chain}; the chain is therefore strictly cheaper and the described strategy is not cost-minimizing.
\end{proof}

\subsection{Proof of Proposition~\ref{prop:fragmentation}}

Recall that Proposition~\ref{prop:fragmentation} relates the cost of the optimal AI deployment strategy for a step sequence to its fragmentation index.
Intuitively, we expect a workflow in which automatable tasks are clustered together will benefit more from AI deployment than one where steps with high and low levels of automation susceptibility are interspersed.
This is due to the potential benefits from AI chains.

First recall the definition of fragmentation index for a given step sequence $\mathcal{S} = (s_1, \dotsc, s_m)$.
As a reminder, we assume for Proposition~\ref{prop:fragmentation} that $\AIskill{i}\AItime{i} = 1$ for all $s_i$.
Suppose that each step $s_i$ succeeds independently with probability $q_i$; any step that does not succeed is said to fail.
Write $F$ for the set of steps that fail, and $\mathcal{C} = \{C_1, \dotsc, C_k\}$ for the random variable representing the collection of maximal connected components of non-failed steps.
The weight of each $C_j \in \mathcal{C}$ is defined to be $\omega(C_j) = \min\{ 1, \sum_{s_i \in C_j} \manualSkill{i}\manualTime{i} \}$.
That is, each $C_j$ has weight $1$ unless the sum of the manual costs of its steps is less than $1$ (due to our assumption that AI management costs are normalized to $1$).

We now introduce some notation: for each $s_i$, write $\minTime{i} = \min\left\{\manualSkill{i}\manualTime{i}, \frac{\AIskill{i}\AItime{i}}{q_i}\right\}$.
That is, $\minTime{i}$ is the minimum cost to implement step $s_i$ individually (whether manually or through AI augmentation).
Given a realization of $\mathcal{C}$ and $F$, we define the realized fragmentation to be
\[ \sum_{s_i \in F}\minTime{i} + \sum_{C_j \in \mathcal{C}} \omega(C_j).\]
The \emph{fragmentation index} is defined to be the expected value of the realized fragmentation.

We now fix the sequence $\mathcal{S} = (s_1, \dotsc, s_m)$, write $FI$ for the fragmentation index of $\mathcal{S}$, and let $OPT$ be the cost of the optimal (i.e., cost-minimizing) task structure for those steps.

We begin by showing that $FI$ is not much larger than $OPT$.

\begin{lem}
\label{lem:FI.upper.bound}
    $FI \leq \tfrac{5}{4}OPT$.
\end{lem}
\begin{proof}
Fix an arbitrary task sequence $\T = (T_1, \dotsc, T_n)$.
Partition this into two sets of tasks: $\T_{\manualLetter}$ containing all (singleton) tasks that are completed manually, and $\T_{\AIletter}$ containing all other tasks (consisting of AI-augmented singletons and AI chains).
We first claim that
\[ FI \leq \sum_{T_j \in \T_{\AIletter}}\left(1 + \sum_{s_i \in T_j}(1-\alpha^{d_i})(1+\minTime{i})\right) + \sum_{(s_i) \in \T_{\manualLetter}}\manualSkill{i}\manualTime{i}. \]

To see why, first note that for any human-executed task $(s_i) \in \T_{\manualLetter}$, either the task fails in the realization of $F$ (in which case it contributes $\minTime{i} \leq \manualSkill{i}\manualTime{i}$ to $FI$), or it succeeds (in which case it appears in some connected component $C_j$, contributing at most $\manualSkill{i}\manualTime{i}$ to the component's weight $\omega(C_j)$).

Next consider the AI-assisted tasks.
For any $T_j \in \T_{\AIletter}$, consider the set of steps in $T_j$ that fail in any given realization of $F$.
If no steps in $T_j$ fail, then $T_j$ collectively contributes at most $1$ to the realized fragmentation (since all $s_i \in T_j$ lie in the same connected component of non-failed tasks).
Otherwise, each failed task $s_i \in T_j$ contributes $\minTime{i}$ to the realized fragmentation (for itself, recalling that $\minTime{i}$ is the minimal cost of executing $s_i$ on its own) plus an additional value of at most $1$ for the weight of a potential new connected component of non-failed tasks.
As each task $s_i \in T_j$ fails independently with probability $(1-\alpha^{d_i})$, we get that the contribution of tasks in $T_j$ to the fragmentation index is at most $1 + \sum_{s_i \in T_j}(1-\alpha^{d_i})(1+\minTime{i})$ as claimed.

On the other hand, recall that if we take $\T$ to be the cost-minimizing task sequence, then
\[ OPT = \sum_{T_j \in \T_{\AIletter}} \alpha^{-d(T_j)} + \sum_{(s_i) \in \T_{\manualLetter}} \manualSkill{i}\manualTime{i} \]
by definition, where we write $d(T_j) = \sum_{s_i \in T_j}d_i$ for the total difficulty of steps in task $T_j$.

Comparing our expression for $OPT$ with our bound on $FI$, we see that each task in $\T_{\manualLetter}$ contributes the same amount to each, whereas each $T_j \in \T_{\AIletter}$ contributes $1+\sum_{s_i \in T_j}(1-\alpha^{d_i})(1+\minTime{i})$ to the former and $\alpha^{-d(T_j)}$ to the latter.
We will show that, for each $T_j \in \T_{\AIletter}$, $1 + \sum_{s_i \in T_j}(1-\alpha^{d_i})(1+\minTime{i}) \leq \frac{5}{4} \alpha^{-d(T_j)}$, which will complete the proof.

First, since $\minTime{i} \leq \alpha^{-d_i}$ for each $s_i \in \mathcal{S}$, we have $1 + \sum_{s_i \in T_j}(1-\alpha^{d_i})(1+\minTime{i}) \leq 1 + \sum_{s_i \in T_j}(1-\alpha^{d_i})(1+\alpha^{-d_i})$.

Next note that for any $x,y\in [0,1]$, $(1-xy)(1+\tfrac{1}{xy}) \geq (1-x)(1+\tfrac{1}{x}) + (1-y)(1+\tfrac{1}{y})$.
(This can be checked by expanding and taking derivatives.)
Repeatedly applying this fact, we get that $1 + \sum_{s_i \in T_j}(1-\alpha^{d_i})(1+\alpha^{-d_i}) \leq 1 + (1-\alpha^{d(T_j)})(1+\alpha^{-d(T_j)}) = 1 + \alpha^{-d(T_j)} - \alpha^{d(T_j)}$ for any $T_j \in \T_{\AIletter}$.
The ratio between this expression and $\alpha^{-d(T_j)}$ is maximized at $\alpha^{-d(T_j)} = 1/2$, achieving a value of $5/4$ as claimed.
\end{proof}

The bound of $\tfrac{5}{4}$ in Lemma~\ref{lem:FI.upper.bound} is tight, as the following example shows.

\begin{example}
\label{ex:FI.tight}
Consider a three-step workflow (recall that $\AIskill{i}\AItime{i} = 1$ for all $i$):
\begin{center}
\begin{tabular}{c|ccc}
\toprule
\makecell{Step\\$(i)$} & \makecell{Manual cost\\$(\manualSkill{i}, \manualTime{i})$} & \makecell{AI management cost\\$(\AIskill{i}, \AItime{i})$} & \makecell{AI success probability\\$(q_i)$} \\
\midrule
1 & $(1, 1)$ & $(1, 1)$ & $1$ \\
2 & $(1, 2)$ & $(1, 1)$ & $1/2$ \\
3 & $(1, 1)$ & $(1, 1)$ & $1$ \\
\bottomrule
\end{tabular}
\end{center}
The first and last steps always succeed, while the second succeeds with probability $1/2$ and has a manual cost of $2$.
The optimal policy automates all three together for an expected cost of $\AIskill{3}\AItime{3}/(q_1 q_2 q_3) = 1/(1 \cdot \tfrac12 \cdot 1) = 2$, whereas the fragmentation index is $\tfrac12 \times 1 + \tfrac12 \times (1+2+1) = \tfrac52$.
The ratio $FI/OPT$ is therefore $\tfrac54$.
\end{example}

We next argue that $FI$ is not much smaller than $OPT$.
We begin with an analysis under the assumption that $\manualSkill{i}\manualTime{i} \geq 1$ for all $s_i \in \mathcal{S}$.
That is, for each step $s_i$, completing the step manually is at least as costly as managing an AI tool to complete the step in a single attempt (with guaranteed success).

\begin{lem}
\label{lem:FI.lower.bound.4}
    Suppose $\manualSkill{i}\manualTime{i} \geq 1$ for all $s_i \in \mathcal{S}$.
    Then $FI \geq \tfrac{1}{4}OPT$.
\end{lem}
\begin{proof}

    Consider the following task sequence.
    Beginning with the first step $s_1$, consider the maximal contiguous sequence of steps $T$ whose success probability $\alpha^{d(T)}$ is greater than or equal to $1/2$.
    If that set is empty (i.e., the first step has success probability strictly less than $1/2$) then the first step is set to be completed individually, either manually or augmented, whichever is cheaper.
    In this case, we'll say the resulting singleton task is completed \emph{individually}.
    Otherwise, the sequence $T = (s_1, \dotsc, s_\ell)$ (which could still be a singleton task) is added to the task sequence as an AI chain; we call this a \emph{non-individual} task.
    This procedure is then repeated beginning with the next step (which is $s_2$ in the former case, or $s_{\ell+1}$ in the latter case), until all steps have been added to a task.
    Call the resulting task sequence $\T$.

    \medskip

    Let $ALG$ denote the cost of the task sequence $\T$.
    Note that we must have $ALG \geq OPT$.
    We will argue that $FI \geq ALG/4$, which will prove the claim.

    \medskip

    We first define some notation.
    Write $\T_I$ for the set of individual tasks in $\T$ and $\T_{NI}$ for the set of non-individual tasks.
    Note that $\alpha^{d_i} < 1/2$ for all $(s_i) \in \T_I$, by construction.
    Note also that $\T_{\manualLetter} \subseteq \T_I$ (recalling that $\T_{\manualLetter}$ is the set of all singleton tasks executed manually), and this containment may be strict if $\T_I$ contains singleton tasks that are augmented.
    We emphasize that $\T_{NI}$ may still include singleton tasks, but any singleton task $(s_i) \in \T_{NI}$ must have the property that $\alpha^{d_i} \geq 1/2$.

    If the final task $s_m$ from $\mathcal{S}$ is contained in some $T_j \in \T_{NI}$, we say that $T_j$ is the \emph{terminal} task.
    All other tasks in $\T_{NI}$ are non-terminal tasks.
Note that if the final step is a singleton task in $\T_I$ then no task is terminal.
    For each non-terminal task $T_j \in \T_{NI}$, we'll write $\overline{T}_j$ to denote $T_j$ plus the first step that follows $T_j$.
    For example, if $T_j = (s_i, \dotsc, s_\ell)$, then $\overline{T}_j = (s_i, \dotsc, s_\ell, s_{\ell+1})$.
    The set $\T_{NI}$ then has the property that for each task $T_j \in \T_{NI}$ we have $\alpha^{d(T_j)} \geq 1/2$, and for each non-terminal $T_j \in \T_{NI}$ we have $\alpha^{(d(\overline{T}_j))} < 1/2$.

    \medskip

    We are now ready to relate $FI$ and $ALG$.
    Let $k = |\T_{NI}|$ be the number of non-individual tasks.
    Note first that
    \[ ALG = \sum_{T_j \in \T_{NI}} \alpha^{-d(T_j)} + \sum_{(s_i) \in \T_I} \minTime{i} \leq 2k + \sum_{(s_i) \in \T_I} \minTime{i}. \]
    This is because $\alpha^{d(T_j)} \geq 1/2$ for each $T_j \in \T_{NI}$, and hence $\alpha^{-d(T_j)} \leq 2$.

    \medskip

    Next we bound $FI$.
    Note that the assumption $\manualSkill{i}\manualTime{i} \geq 1$ implies that, in any realization of the connected components of non-failed steps $\mathcal{C} = (C_1, \dotsc, C_\ell)$, $\omega(C_j) = 1$ for all $j$.
    This implies that the realized fragmentation is equal to $1$ plus, for each $s_i \in F$ (i.e., each step $s_i$ that fails), a contribution of $\minTime{i}$ plus an additional $1$ in the event that $i > 1$ and the task immediately preceding $s_i$ did not fail.
    (This additional cost of $1$ accounts for creating a new connected component of non-failed steps ending with $s_{i-1}$.)

    We claim that $FI \geq k/2 + \frac{1}{2}\sum_{(s_i) \in \T_I} \minTime{i}$.
    We will show this by charging  the realized fragmentation to tasks in $\T_{NI}$ and $\T_I$, so that each $T_j \in \T_{NI}$ is charged at least $1$ with probability at least $1/2$, and each $(s_i) \in \T_I$ is charged at least $\minTime{i}$ with probability at least $1/2$.

    To see this, let $E_j$ denote the event that some step in $\overline{T}_j$ fails, for each non-terminal $T_j \in \T_{NI}$.
    Note that $E_j$ occurs with probability at least $1/2$, since $\alpha^{d(\overline{T}_j)} < 1/2$.
    Also, if $E_j$ occurs, then either a step $s_i \in T_j$ fails, or no step in $T_j$ fails but the step immediately following $T_j$ does.
    In the former case, we will charge the $\minTime{i} \geq 1$ contribution of $s_i$'s failure to $T_j$.
    In the latter case, it must be that the step immediately preceding the failed task did not fail; so we will charge the additional contribution of $1$ from $s_i$'s failure (as described above) to $T_j$.
    Additionally, if there is a terminal $T_j$ in $\T_{NI}$, then we charge the baseline $1$ from the calculation of $FI$ to that terminal $T_j$.
    Aggregating over all these cases, we have that at least $1$ is charged to each set $T_j \in \T_{NI}$ with probability at least $1/2$.
    Finally, for each $(s_i) \in \T_I$ that fails (which occurs with probability at least $1/2$, by construction), we charge the $\minTime{i}$ portion of its contribution to task $(s_i)$.
    We conclude that $FI \geq k/2 + \frac{1}{2}\sum_{(s_i) \in \T_I}\minTime{i}$ as claimed.

    But now since $FI \geq k/2 + \frac{1}{2}\sum_{(s_i) \in \T_I}\minTime{i}$ and $ALG \leq 2k + \sum_{(s_i) \in \T_I}\minTime{i}$, we conclude $FI \geq ALG/4$ as claimed.
\end{proof}

The approximation factor in Lemma~\ref{lem:FI.lower.bound.4} cannot be improved below $\tfrac{2}{3}(2+\sqrt{2}) \approx 2.276$, although we do not have a matching lower bound of $4$.

\begin{example}
\label{ex:FI.lower.gap}
Consider a sequence of $m$ identical steps:
\begin{center}
\begin{tabular}{c|ccc}
\toprule
\makecell{Step\\$(i)$} & \makecell{Manual cost\\$(\manualSkill{i}, \manualTime{i})$} & \makecell{AI management cost\\$(\AIskill{i}, \AItime{i})$} & \makecell{AI success probability\\$(q_i)$} \\
\midrule
All $i$ & $(1, \sqrt{2})$ & $(1, 1)$ & $1/\sqrt{2}$ \\
\bottomrule
\end{tabular}
\end{center}
The task sequence described in the proof of Lemma~\ref{lem:FI.lower.bound.4} handles each task separately, for a total cost of $m \sqrt{2}$, so the optimal task sequence performs at least this well.
The fragmentation index is such that each task contributes $\sqrt{2}$ if it fails, plus $1$ if the preceding task did not fail, for a total contribution of $(1-1/\sqrt{2})(\sqrt{2} + 1/\sqrt{2}) = \tfrac{3}{2}(\sqrt{2}-1) \approx 0.6213$.
As $m$ grows large, the ratio between $m\sqrt{2}$ and $1 + m \times 0.6213$ approaches $\tfrac{2}{3}(2+\sqrt{2}) \approx 2.276$.
\end{example}

We can relax the assumption in Lemma~\ref{lem:FI.lower.bound.4} that $\manualSkill{i}\manualTime{i} \geq 1$ for all $s_i$, but at the cost of a worse approximation factor.

\begin{lem}
\label{lem:FI.lower.bound.8}
    For general $\manualSkill{i}\manualTime{i}$ (i.e., allowing $\manualSkill{i}\manualTime{i} < 1$ for some steps $s_i$), $FI \geq \frac{1}{8}OPT$.
\end{lem}
\begin{proof}
    The argument follows the proof of Lemma~\ref{lem:FI.lower.bound.4}, with the following changes.

    First we make a slight change in the definition of the task sequence $\T$ that defines $ALG$.
    If a constructed task $T_j$ has the property that $\sum_{s_i \in T_j} \manualSkill{i}\manualTime{i} < 1$ (which implies it is cheaper to run all steps of $T_j$ manually than as an AI chain that is guaranteed to succeed), we switch to running the tasks of $T_j$ manually in the task sequence.
    In our analysis, we still consider $T_j$ to be a set in $\T_{NI}$; but its cost is taken to be $\sum_{s_i \in T_j} \manualSkill{i}\manualTime{i}$.

    This modification changes our upper bound on the total cost of $ALG$ to the following:
    \[ ALG \leq 2 \sum_{T_j \in \T_{NI}} \min\left\{1, \sum_{s_i \in T_j} \manualSkill{i}\manualTime{i} \right\} + \sum_{ (s_i) \in \T_I} \minTime{i}. \]

    Then, to bound $FI$, we claim that
    \[ FI \geq \frac{1}{4}\sum_{T_j \in \T_{NI}} \min\left\{1, \sum_{s_i \in T_j} \manualSkill{i}\manualTime{i} \right\} + \frac{1}{2}\sum_{(s_i) \in \T_I} \minTime{i} \]
    which would prove the claim.

    To see why this bound on $FI$ holds, we employ a charging argument as before.
    Recall that for each non-terminal $T_j \in \T_{NI}$, the event $E_j$ occurs with probability at least $1/2$.
    When it does, we charge to $T_j$ the contribution $\minTime{i}$ to $FI$ from any $s_i \in F$ that lies in $T_j$.
    Furthermore, when $E_j$ occurs, for each connected component $C_\ell$ that intersects $T_j$ we additionally charge the contribution $\omega(C_\ell)$ from $FI$ to $T_j$.
    Crucially, the contribution from each $C_\ell$ can be charged to at most two different tasks $T_j$ in this way: once for the leftmost $T_j$ that it intersects, and once for the rightmost $T_j$ that it intersects.
    (The reason is that if some $T_j$ is a subset of $C_\ell$, then by definition no step of $T_j$ fails and hence $E_j$ did not occur.
    So this charging can only occur for a task $T_j$ that intersects $C_\ell$ but is not a subset of $C_\ell$, of which there are at most two.)
    Moreover, this charging to $T_j$ adds up to a  total value of at least $\min\left\{1, \sum_{s_i \in T_j} \manualSkill{i}\manualTime{i} \right\}$.
    Since the charging occurs with probability $1/2$ for each $T_j \in \T_{NI}$, and we are at most double-counting each $\omega(C_\ell)$, we conclude that the expected sum of all $\omega(C_\ell)$, plus $\minTime{i}$ for all failed tasks $s_i$ that lie in some $T_j \in \T_{NI}$, is at least $1/4$ of $\sum_{T_j \in \T_{NI}}\min\left\{1, \sum_{s_i \in T_j} \manualSkill{i}\manualTime{i} \right\}$.

    The charging for $(s_i) \in \T_I$ remains unchanged: the value $\minTime{i}$ is charged in the event that step $s_i$ fails, which occurs with probability at least $1/2$.

    Taken together, we obtain our desired $8$ approximation.
\end{proof}

We suspect the factor of $8$ in Lemma~\ref{lem:FI.lower.bound.8} is not tight.
But, as the following example shows, the factor is strictly greater than the factor $4$ from Lemma~\ref{lem:FI.lower.bound.4}, so the assumption that $\manualSkill{i}\manualTime{i} \geq 1$ for all $s_i$ is necessary for Lemma~\ref{lem:FI.lower.bound.4} to hold.

\begin{example}
\label{ex:FI.necessity}
Consider a sequence that alternates between two step types, with $K > 1$ such pairs (so $2K$ steps in total), where $\epsilon$ is vanishingly small:
\begin{center}
\begin{tabular}{c|ccc}
\toprule
\makecell{Step\\$(i)$} & \makecell{Manual cost\\$(\manualSkill{i}, \manualTime{i})$} & \makecell{AI management cost\\$(\AIskill{i}, \AItime{i})$} & \makecell{AI success probability\\$(q_i)$} \\
\midrule
Odd $i$ & $(1, 0)$ & $(1, 1)$ & $1/\sqrt{2}-\epsilon$ \\
Even $i$ & $(1, 1)$ & $(1, 1)$ & $1$ \\
\bottomrule
\end{tabular}
\end{center}
The task sequence described in the proof of Lemma~\ref{lem:FI.lower.bound.8} groups the steps into pairs of two-step tasks, for a total cost of $K \sqrt{2}$ (ignoring $\epsilon$).
The fragmentation index is $1$ plus $1$ for each task that fails, which is $1 + K(1 - 1/\sqrt{2})$ (again ignoring the impact of $\epsilon$).
As $K$ grows large, the ratio tends to $2(\sqrt{2}+1)$, which is strictly greater than $4$.
\end{example}

\begin{proof}[Proof of Proposition~\ref{prop:fragmentation}]
Combining Lemmas~\ref{lem:FI.upper.bound}, \ref{lem:FI.lower.bound.4}, and~\ref{lem:FI.lower.bound.8} yields Proposition~\ref{prop:fragmentation}: in general $\tfrac{1}{8}OPT \leq FI \leq \tfrac{5}{4}OPT$, and under the additional assumption $\manualSkill{i}\manualTime{i} \geq 1$ for all $i$, $\tfrac{1}{4}OPT \leq FI \leq \tfrac{5}{4}OPT$.
\end{proof}

\subsection{Proof of Proposition~\ref{prop:reorg_threshold}}

\begin{proof}
Each $C_\T$ is a constant plus terms $\AIskill{r_c}\AItime{r_c}\,\alpha^{-D_c}$ with $D_c \ge 1$, with derivatives
\[
\begin{aligned}
C_\T'(\alpha) &= -\sum_{c \in \mathcal{C}(\T)} \AIskill{r_c}\AItime{r_c}\,D_c\,\alpha^{-D_c-1} \;\le\; 0, \\
C_\T''(\alpha) &= \sum_{c \in \mathcal{C}(\T)} \AIskill{r_c}\AItime{r_c}\,D_c(D_c+1)\,\alpha^{-D_c-2} \;\ge\; 0,
\end{aligned}
\]
so each $C_\T$ is non-increasing and convex.

Within a regime where $\T$ is optimal, $g^*(\alpha) = g_\T(\alpha) = -C_\T'(\alpha) \ge 0$, and the convexity of $C_\T$ makes it non-increasing in $\alpha$ (when $\T$ uses no chain $g^*(\alpha) = 0$ but is otherwise strictly decreasing).
The returns to a fixed configuration in a regime thus weakly diminish.
This proves the first claim of the proposition.

Next, let $\alpha_0 \in (0,1)$ be a threshold where the optimum switches from $\T$ (optimal just below $\alpha_0$) to $\T'$ (just above $\alpha_0$).\footnote{We assume this threshold exists. The only scenario where no such threshold exists is the case in which all steps are executed manually for all $\alpha \in (0,1)$. This is the trivial, uninteresting case and is assumed away.}
Define
\[
\phi(\alpha) = C_{\T'}(\alpha) - C_\T(\alpha).
\]
The function $\phi(\alpha)$ is the difference of two finite sums of powers of $\alpha$ and hence continuous and differentiable on $(0,1)$. 
The two costs are equal at $\alpha_0$, with $C_{\T'}$ the cheaper just above and $C_\T$ the cheaper just below.
Therefore, by definition $\phi$ satisfies:
\[
\phi(\alpha_0^{-}) \ge 0,
\qquad
\phi(\alpha_0) = 0,
\qquad
\phi(\alpha_0^{+}) \le 0.
\]
A differentiable function with such a sign change has nonpositive derivative in the local neighborhood of $\alpha_0$.
So $\phi'(\alpha_0) \le 0$, or equivalently, $C_{\T'}'(\alpha_0) \le C_\T'(\alpha_0)$.
By the definition $g_\T(\alpha) = -\,\mathrm{d}C_\T/\mathrm{d}\alpha$, this is
\[
g_\T(\alpha_0) \;=\; -C_\T'(\alpha_0)
\;\le\;
-C_{\T'}'(\alpha_0) \;=\; g_{\T'}(\alpha_0).
\]
Since $g^*(\alpha)$ coincides with $g_\T(\alpha)$ just below $\alpha_0$ and with $g_{\T'}(\alpha)$ just above, the firm's marginal benefit function $g^*(\alpha)$ jumps upward as $\alpha$ crosses $\alpha_0$.
\end{proof}
